% GENERATED FILE. Edit canonical lessons, then run npm run generate:books.
% canonical-source-hash: fnv1a64:1bfd1984f69bfc1b
% canonical-lessons: ML-C548-mekkanikku, ML-C548-vyapari, ML-C548-matsyattolilali, ML-C548-pacakakkaran, ML-C548-phottographar, ML-R548-review-541-548

\chapter{Mechanics, Merchants and Cooks}
\label{ch:ml-mechanics-merchants-and-cooks}
\hlchaptermodality{548}

\begin{chapteropening}
\textbf{By the end of this chapter:} \emph{I can say a mechanic, a merchant, a fisherman, a professional cook and a photographer.}

Complete the last lesson of chapter 548: I can say a mechanic, a merchant, a fisherman, a professional cook and a photographer.
\end{chapteropening}

\section[\texorpdfstring{\ml{മെക്കാനിക്ക്} (mekkānikkŭ)}{മെക്കാനിക്ക് (mekkānikkŭ)}]{\ml{മെക്കാനിക്ക്} (mekkānikkŭ) — a mechanic}

\label{lesson:ML-C548-mekkanikku}



\begin{quote}
Before the new one: say the Malayalam for an electrician, then the Malayalam for a plumber.
\end{quote}

\subsection*{You'll want to know: \ml{മെക്കാനിക്ക്}}
\textbf{\ml{മെക്കാനിക്ക്}} — \emph{mekkānikkŭ} — \textquotedblleft{}a mechanic\textquotedblright{}.

\begin{quote}
\ml{വണ്ടി} \ml{കേടായപ്പോൾ} \ml{മെക്കാനിക്ക്} \ml{വന്നു}.
\end{quote}

\begin{quote}
\emph{vaṇṭi kēṭāyappōḷ mekkānikkŭ vannu.}
\end{quote}

\begin{quote}
When the vehicle broke down, the mechanic came.
\end{quote}

A garage is a \ml{വർക്ക്ഷോപ്പ്} in everyday talk.

\begin{cousinweb}
\textbf{\ml{മെക്കാനിക്ക്}} (\emph{mekkānikkŭ}) was borrowed from English \emph{mechanic}.
\end{cousinweb}

\begin{grammarlens}[title={What you've built}]
One more word for describing what you lived and felt.
\end{grammarlens}

\subsection*{Your turn}
\begin{itemize}
\raggedright
  \item \emph{Say it:} \emph{mekkānikkŭ}
  \item \emph{Say it:} \emph{mekkānikkŭ}, once more
  \item \emph{Say it:} \emph{vaṇṭi kēṭāyappōḷ mekkānikkŭ vannu.} — When the vehicle broke down, the mechanic came.
  \item \emph{Recall:} say the Malayalam for an electrician, then the Malayalam for a plumber, then say \emph{mekkānikkŭ} again
\end{itemize}

\begin{tcolorbox}[breakable,colback=teal!4,colframe=teal!35!black,title={Before you move on}]
What does \ml{മെക്കാനിക്ക്} mean? (\textquotedblleft{}A mechanic\textquotedblright{}.) Say it once more.
\end{tcolorbox}
\section[\texorpdfstring{\ml{വ്യാപാരി} (vyāpāri)}{വ്യാപാരി (vyāpāri)}]{\ml{വ്യാപാരി} (vyāpāri) — a merchant; a dealer}

\label{lesson:ML-C548-vyapari}



\begin{quote}
Before the new one: say the Malayalam for a plumber, then the Malayalam for a mechanic.
\end{quote}

\subsection*{You'll want to know: \ml{വ്യാപാരി}}
\textbf{\ml{വ്യാപാരി}} — \emph{vyāpāri} — \textquotedblleft{}a merchant; a dealer\textquotedblright{}.

\begin{quote}
\ml{ആ} \ml{വ്യാപാരി} \ml{ചന്തയിൽ} \ml{അരി} \ml{വിൽക്കുന്നു}.
\end{quote}

\begin{quote}
\emph{ā vyāpāri cantayil ari vilkkunnu.}
\end{quote}

\begin{quote}
That merchant sells rice at the market.
\end{quote}

\ml{വ്യാപാരം} is trade; a bigger dealer than a \ml{കച്ചവടക്കാരൻ}.

\begin{cousinweb}
\textbf{\ml{വ്യാപാരി}} (\emph{vyāpāri}) is built from Sanskrit \emph{vyāpāra}, \textquotedblleft{}occupation, business\textquotedblright{} (CDIAL 12205), with the ending -\ml{ഇ} for a person.
\end{cousinweb}

\begin{grammarlens}[title={What you've built}]
One more word for describing what you lived and felt.
\end{grammarlens}

\subsection*{Your turn}
\begin{itemize}
\raggedright
  \item \emph{Say it:} \emph{vyāpāri}
  \item \emph{Say it:} \emph{vyāpāri}, once more
  \item \emph{Say it:} \emph{ā vyāpāri cantayil ari vilkkunnu.} — That merchant sells rice at the market.
  \item \emph{Recall:} say the Malayalam for a plumber, then the Malayalam for a mechanic, then say \emph{vyāpāri} again
\end{itemize}

\begin{tcolorbox}[breakable,colback=teal!4,colframe=teal!35!black,title={Before you move on}]
What does \ml{വ്യാപാരി} mean? (\textquotedblleft{}A merchant; a dealer\textquotedblright{}.) Say it once more.
\end{tcolorbox}
\section[\texorpdfstring{\ml{മത്സ്യത്തൊഴിലാളി} (matsyattoḻilāḷi)}{മത്സ്യത്തൊഴിലാളി (matsyattoḻilāḷi)}]{\ml{മത്സ്യത്തൊഴിലാളി} (matsyattoḻilāḷi) — a fisherman; a fish worker}

\label{lesson:ML-C548-matsyattolilali}



\begin{quote}
Before the new one: say the Malayalam for a mechanic, then the Malayalam for a merchant.
\end{quote}

\subsection*{You'll want to know: \ml{മത്സ്യത്തൊഴിലാളി}}
\textbf{\ml{മത്സ്യത്തൊഴിലാളി}} — \emph{matsyattoḻilāḷi} — \textquotedblleft{}a fisherman; a fish worker\textquotedblright{}.

\begin{quote}
\ml{മത്സ്യത്തൊഴിലാളികൾ} \ml{പുലർച്ചെ} \ml{കടലിൽ} \ml{പോയി}.
\end{quote}

\begin{quote}
\emph{matsyattoḻilāḷikaḷ pularcce kaṭalil pōyi.}
\end{quote}

\begin{quote}
The fishermen went out to sea at dawn.
\end{quote}

Plural: \ml{മത്സ്യത്തൊഴിലാളികൾ}, \emph{matsyattoḻilāḷikaḷ}.

\ml{മത്സ്യം} fish + \ml{തൊഴിലാളി} a worker; the formal word, as in the news.

\begin{cousinweb}
\textbf{\ml{മത്സ്യത്തൊഴിലാളി}} (\emph{matsyattoḻilāḷi}) is a Malayalam compound built from Sanskrit \emph{matsya}, \textquotedblleft{}fish\textquotedblright{} (CDIAL 9758), and \ml{തൊഴിലാളി}, \textquotedblleft{}a worker\textquotedblright{}, which is native (DEDR 3524).
\end{cousinweb}

\begin{grammarlens}[title={What you've built}]
One more word for describing what you lived and felt.
\end{grammarlens}

\subsection*{Your turn}
\begin{itemize}
\raggedright
  \item \emph{Say it:} \emph{matsyattoḻilāḷi}
  \item \emph{Say it:} \emph{matsyattoḻilāḷi}, once more
  \item \emph{Say it:} \emph{matsyattoḻilāḷikaḷ pularcce kaṭalil pōyi.} — The fishermen went out to sea at dawn.
  \item \emph{Recall:} say the Malayalam for a mechanic, then the Malayalam for a merchant, then say \emph{matsyattoḻilāḷi} again
\end{itemize}

\begin{tcolorbox}[breakable,colback=teal!4,colframe=teal!35!black,title={Before you move on}]
What does \ml{മത്സ്യത്തൊഴിലാളി} mean? (\textquotedblleft{}A fisherman; a fish worker\textquotedblright{}.) Say it once more.
\end{tcolorbox}
\section[\texorpdfstring{\ml{പാചകക്കാരൻ} (pācakakkāran)}{പാചകക്കാരൻ (pācakakkāran)}]{\ml{പാചകക്കാരൻ} (pācakakkāran) — a professional cook}

\label{lesson:ML-C548-pacakakkaran}



\begin{quote}
Before the new one: say the Malayalam for a merchant, then the Malayalam for a fisherman.
\end{quote}

\subsection*{You'll want to know: \ml{പാചകക്കാരൻ}}
\textbf{\ml{പാചകക്കാരൻ}} — \emph{pācakakkāran} — \textquotedblleft{}a professional cook\textquotedblright{}.

\begin{quote}
\ml{ഹോട്ടലിലെ} \ml{പാചകക്കാരൻ} \ml{നല്ല} \ml{ബിരിയാണി} \ml{ഉണ്ടാക്കി}.
\end{quote}

\begin{quote}
\emph{hōṭṭalile pācakakkāran nalla biriyāṇi uṇṭākki.}
\end{quote}

\begin{quote}
The cook at the restaurant made a good biryani.
\end{quote}

\ml{പാചകം} cooking + -\ml{ക്കാരൻ}, one whose work it is.

\begin{cousinweb}
\textbf{\ml{പാചകക്കാരൻ}} (\emph{pācakakkāran}) is built from \ml{പാചകം}, from Sanskrit \emph{pācaka}, \textquotedblleft{}cooking, roasting\textquotedblright{} (Monier-Williams), with the Malayalam ending -\ml{ക്കാരൻ}, \textquotedblleft{}one who does it\textquotedblright{}.
\end{cousinweb}

\begin{grammarlens}[title={What you've built}]
One more word for describing what you lived and felt.
\end{grammarlens}

\subsection*{Your turn}
\begin{itemize}
\raggedright
  \item \emph{Say it:} \emph{pācakakkāran}
  \item \emph{Say it:} \emph{pācakakkāran}, once more
  \item \emph{Say it:} \emph{hōṭṭalile pācakakkāran nalla biriyāṇi uṇṭākki.} — The cook at the restaurant made a good biryani.
  \item \emph{Recall:} say the Malayalam for a merchant, then the Malayalam for a fisherman, then say \emph{pācakakkāran} again
\end{itemize}

\begin{tcolorbox}[breakable,colback=teal!4,colframe=teal!35!black,title={Before you move on}]
What does \ml{പാചകക്കാരൻ} mean? (\textquotedblleft{}A professional cook\textquotedblright{}.) Say it once more.
\end{tcolorbox}
\section[\texorpdfstring{\ml{ഫോട്ടോഗ്രാഫർ} (phōṭṭōgrāphar)}{ഫോട്ടോഗ്രാഫർ (phōṭṭōgrāphar)}]{\ml{ഫോട്ടോഗ്രാഫർ} (phōṭṭōgrāphar) — a photographer}

\label{lesson:ML-C548-phottographar}



\begin{quote}
Before the new one: say the Malayalam for a fisherman, then the Malayalam for a professional cook.
\end{quote}

\subsection*{You'll want to know: \ml{ഫോട്ടോഗ്രാഫർ}}
\textbf{\ml{ഫോട്ടോഗ്രാഫർ}} — \emph{phōṭṭōgrāphar} — \textquotedblleft{}a photographer\textquotedblright{}.

\begin{quote}
\ml{കല്യാണത്തിന്} \ml{രണ്ട്} \ml{ഫോട്ടോഗ്രാഫർമാർ} \ml{ഉണ്ടായിരുന്നു}.
\end{quote}

\begin{quote}
\emph{kalyāṇattinŭ raṇṭŭ phōṭṭōgrāpharmār uṇṭāyirunnu.}
\end{quote}

\begin{quote}
There were two photographers at the wedding.
\end{quote}

Plural: \ml{ഫോട്ടോഗ്രാഫർമാർ}, \emph{phōṭṭōgrāpharmār}.

Plural: \ml{ഫോട്ടോഗ്രാഫർമാർ}, with -\ml{മാർ} for people.

\begin{cousinweb}
\textbf{\ml{ഫോട്ടോഗ്രാഫർ}} (\emph{phōṭṭōgrāphar}) was borrowed from English \emph{photographer}.
\end{cousinweb}

\begin{grammarlens}[title={What you've built}]
One more word for describing what you lived and felt.
\end{grammarlens}

\subsection*{Your turn}
\begin{itemize}
\raggedright
  \item \emph{Say it:} \emph{phōṭṭōgrāphar}
  \item \emph{Say it:} \emph{phōṭṭōgrāphar}, once more
  \item \emph{Say it:} \emph{kalyāṇattinŭ raṇṭŭ phōṭṭōgrāpharmār uṇṭāyirunnu.} — There were two photographers at the wedding.
  \item \emph{Recall:} say the Malayalam for a fisherman, then the Malayalam for a professional cook, then say \emph{phōṭṭōgrāphar} again
\end{itemize}

\begin{tcolorbox}[breakable,colback=teal!4,colframe=teal!35!black,title={Before you move on}]
What does \ml{ഫോട്ടോഗ്രാഫർ} mean? (\textquotedblleft{}A photographer\textquotedblright{}.) Say it once more.
\end{tcolorbox}
\section[(dialogue)]{Words from chapters 541–548 — a review}

\label{lesson:ML-R548-review-541-548}



\begin{quote}
Say the Malayalam for to subside and to get worse.
\end{quote}

\subsection*{Your turn}
Say these aloud:

\begin{itemize}
\raggedright
  \item to subside, to get worse, to spread by contact, to go without food, to be digested
  \item a trade, a salaried job, a promotion, pay, to step down
  \item a bonus, to apply, to take on staff, to dismiss, an employer
  \item overtime, a venture, an entrepreneur, a work shift, to be interrupted
  \item a superior, being in charge, a job transfer, planning, to cooperate
  \item manual labour, a livelihood, a workload, work experience, to lead
  \item a pilot, a dentist, an architect, an electrician, a plumber
  \item a mechanic, a merchant, a fisherman, a professional cook, a photographer
\end{itemize}

\begin{tcolorbox}[breakable,colback=teal!4,colframe=teal!35!black,title={Before you move on}]
To subside? (\textbf{\ml{ശമിക്കുക}}, \emph{śamikkuka}.) To get worse? (\textbf{\ml{വഷളാകുക}}, \emph{vaṣaḷākuka}.) To spread by contact? (\textbf{\ml{പകരുക}}, \emph{pakaruka}.) To go without food? (\textbf{\ml{ഉപവസിക്കുക}}, \emph{upavasikkuka}.) To be digested? (\textbf{\ml{ദഹിക്കുക}}, \emph{dahikkuka}.) A trade? (\textbf{\ml{തൊഴിൽ}}, \emph{toḻil}.) A salaried job? (\textbf{\ml{ഉദ്യോഗം}}, \emph{udyōgaṁ}.) A promotion? (\textbf{\ml{സ്ഥാനക്കയറ്റം}}, \emph{sthānakkayaṟṟaṁ}.) Pay? (\textbf{\ml{വേതനം}}, \emph{vētanaṁ}.) To step down? (\textbf{\ml{സ്ഥാനമൊഴിയുക}}, \emph{sthānamoḻiyuka}.) A bonus? (\textbf{\ml{ബോണസ്}}, \emph{bōṇasŭ}.) To apply? (\textbf{\ml{അപേക്ഷിക്കുക}}, \emph{apēkṣikkuka}.) To take on staff? (\textbf{\ml{ജോലിക്കെടുക്കുക}}, \emph{jōlikkeṭukkuka}.) To dismiss? (\textbf{\ml{പിരിച്ചുവിടുക}}, \emph{piriccuviṭuka}.) An employer? (\textbf{\ml{തൊഴിലുടമ}}, \emph{toḻiluṭama}.) Overtime? (\textbf{\ml{ഓവർടൈം}}, \emph{ōvarṭaiṁ}.) A venture? (\textbf{\ml{സംരംഭം}}, \emph{saṁrambhaṁ}.) An entrepreneur? (\textbf{\ml{സംരംഭകൻ}}, \emph{saṁrambhakan}.) A work shift? (\textbf{\ml{ഷിഫ്റ്റ്}}, \emph{ṣiphṟṟŭ}.) To be interrupted? (\textbf{\ml{മുടങ്ങുക}}, \emph{muṭaṅṅuka}.) A superior? (\textbf{\ml{മേലുദ്യോഗസ്ഥൻ}}, \emph{mēludyōgasthan}.) Being in charge? (\textbf{\ml{ചുമതല}}, \emph{cumatala}.) A job transfer? (\textbf{\ml{സ്ഥലംമാറ്റം}}, \emph{sthalaṁmāṟṟaṁ}.) Planning? (\textbf{\ml{ആസൂത്രണം}}, \emph{āsūtraṇaṁ}.) To cooperate? (\textbf{\ml{സഹകരിക്കുക}}, \emph{sahakarikkuka}.) Manual labour? (\textbf{\ml{കൂലിപ്പണി}}, \emph{kūlippaṇi}.) A livelihood? (\textbf{\ml{ഉപജീവനം}}, \emph{upajīvanaṁ}.) A workload? (\textbf{\ml{ജോലിഭാരം}}, \emph{jōlibhāraṁ}.) Work experience? (\textbf{\ml{പ്രവൃത്തിപരിചയം}}, \emph{pravṛttiparicayaṁ}.) To lead? (\textbf{\ml{നയിക്കുക}}, \emph{nayikkuka}.) A pilot? (\textbf{\ml{പൈലറ്റ്}}, \emph{pailaṟṟŭ}.) A dentist? (\textbf{\ml{ദന്തവൈദ്യൻ}}, \emph{dantavaidyan}.) An architect? (\textbf{\ml{വാസ്തുശില്പി}}, \emph{vāstuśilpi}.) An electrician? (\textbf{\ml{ഇലക്ട്രീഷ്യൻ}}, \emph{ilakṭrīṣyan}.) A plumber? (\textbf{\ml{പ്ലംബർ}}, \emph{plambar}.) A mechanic? (\textbf{\ml{മെക്കാനിക്ക്}}, \emph{mekkānikkŭ}.) A merchant? (\textbf{\ml{വ്യാപാരി}}, \emph{vyāpāri}.) A fisherman? (\textbf{\ml{മത്സ്യത്തൊഴിലാളി}}, \emph{matsyattoḻilāḷi}.) A professional cook? (\textbf{\ml{പാചകക്കാരൻ}}, \emph{pācakakkāran}.) A photographer? (\textbf{\ml{ഫോട്ടോഗ്രാഫർ}}, \emph{phōṭṭōgrāphar}.)
\end{tcolorbox}
